%----------------------------------------------------------------------
\section{Model computations and design justification}
%----------------------------------------------------------------------
\textbf{Notation.} Nodes $\mathcal{V}=\{\mathrm{Res},\mathrm{T1},\mathrm{U},\mathrm{S1},\mathrm{S2},\mathrm{S3},\mathrm{S4}\}$;
rain gauges $\mathcal{G}=\{\mathrm{R1},\mathrm{R2}\}$; stations $\mathcal{S}=\{\mathrm{S1},\dots,\mathrm{S4}\}$.
$\mathcal{N}^{\uparrow}(i)$ is the set of sources with an edge into $i$ (Figure~\ref{fig:graph}).
$x_{j,t}$ is the aligned feature vector of source $j$ at grid step $t$ (Table~\ref{tab:data}).
$y_i(t)$ is the level of station $i$, $\tau^{+}=\tau+6\Delta$, and lags $\ell$ are counted in
10-minute steps \emph{back from the target time} $\tau^{+}$.

\subsection*{Main blocks}
\textbf{(1) Per-source causal encoders} (learned). A causal dilated temporal convolution (TCN:
kernel 3, dilations $1,2,\dots,32$, receptive field 127 steps $\ge L$), a source-type-specific
projection to width $d$, and a learned node embedding $\mathbf{p}_j$:
\begin{equation}
\mathbf{z}_{j,t}=W^{(s)}_{\mathrm{proj}}\,\mathrm{TCN}^{(s)}\!\big(x_{j,\le t}\big)+\mathbf{p}_j\in\mathbb{R}^{d},\qquad t\le\tau ,
\end{equation}
where $s\in\{\text{level},\text{rain},\text{release}\}$ and the mask and age are input channels.
Causality guarantees that $\mathbf{z}_{j,t}$ never sees data after $t$. The reservoir sequence is
extended into the future with the announced schedule:
$\mathbf{z}_{\mathrm{Res},\tau+h}=\mathrm{MLP}\big([Q^{\mathrm{sch}}_{\tau+h},a_{\tau+h},h]\big)$ for $h=1,\dots,6$.
The latent node U has no sensor. Its sequence is a learned rainfall-runoff transform of both
gauges, $\mathbf{z}_{\mathrm{U},t}=\mathrm{TCN}^{(\mathrm{U})}\big(x_{\mathrm{R1},\le t},x_{\mathrm{R2},\le t}\big)+\mathbf{p}_{\mathrm{U}}$,
and it is trained only through its effect on S3 and S4.

\textbf{(2a) Delay-attention layer} (learned weights on a physically restricted support). For an
edge $j\to i$, the admissible lags are those in the travel-time window $\Lambda_{ij}$ (A8) that refer
to known data:
\begin{equation}
\mathcal{A}_{ij}=\big\{\ell\in\Lambda_{ij}:\ \tau^{+}-\ell\Delta\le\tau\big\}\ \ (\ell\ge6)\ \text{for observed sources},\qquad
\mathcal{A}_{\mathrm{Res},i}=\Lambda_{\mathrm{Res},i}\ \text{(schedule known)} .
\end{equation}
If $\Lambda_{ij}$ lies entirely below 6, i.e.\ the travel time is under one hour, then
$\mathcal{A}_{ij}=\{6,\dots,6+3\}$: the most recent upstream state, whose TCN encoding contains
its trend. The message is a soft choice of the arrival lag:
\begin{align}
\alpha_{ij,\ell}&=\operatorname*{softmax}_{\ell\in\mathcal{A}_{ij}}
\frac{\big(W_Q\mathbf{z}_{i,\tau}\big)^{\!\top}W_K\big[\mathbf{z}_{j,\tau^{+}-\ell}\,\Vert\,\mathbf{e}_{ij}\,\Vert\,\phi(\ell)\big]}{\sqrt d},
&\mathbf{m}_{ij}&=\sum_{\ell\in\mathcal{A}_{ij}}\alpha_{ij,\ell}\,W_V\mathbf{z}_{j,\tau^{+}-\ell},\\
\beta_{ij}&=\operatorname*{softmax}_{j\in\mathcal{N}^{\uparrow}(i)}\mathbf{w}_\beta^{\!\top}\tanh\!\big(W_\beta[\mathbf{z}_{i,\tau}\Vert\mathbf{m}_{ij}\Vert\mathbf{e}_{ij}]\big),
&\mathbf{h}^{(1)}_{i}&=\mathrm{LN}\Big(\mathbf{z}_{i,\tau}+\mathrm{MLP}\big(\textstyle\sum_{j\in\mathcal{N}^{\uparrow}(i)}\beta_{ij}\mathbf{m}_{ij}\big)\Big),
\end{align}
with $\phi(\ell)$ a learned lag embedding. The query $\mathbf{z}_{i,\tau}$ encodes the current
state of the receiving station (its level and trend), so in high flow the attention moves to
shorter lags. \emph{Changing travel times} are thus a learned mapping, and only the window
$\Lambda_{ij}$ is a physical assumption. The weights $\beta_{ij}$ decide how much each upstream source
contributes (for example T1 versus S1 for S2).

\textbf{(2b) Downstream message passing on node states} (learned). One more layer on the
current node states,
$\mathbf{h}^{(2)}_{i}=\mathrm{LN}\big(\mathbf{h}^{(1)}_{i}+\mathrm{MLP}\big(\sum_{j\in\mathcal{N}^{\uparrow}(i)}\gamma_{ij}W_M\mathbf{h}^{(1)}_{j}\big)\big)$
with attention weights $\gamma_{ij}$ computed like $\beta_{ij}$. It lets a station use what
upstream nodes have already worked out about themselves. For example, S3 receives the summary of U
(rain-driven inflow) and of S2 (its own upstream context).

\textbf{(3) Fusion and forecast mapping.} For each station $i\in\mathcal{S}$, a GRU over its own
encoded history gives $\mathbf{s}_i$. A gate decides how far to trust the graph context:
\begin{equation}
\mathbf{g}_i=\sigma\big(W_g[\mathbf{s}_i\Vert\mathbf{h}^{(2)}_i]\big),\qquad
\mathbf{f}_i=\mathbf{g}_i\odot\mathbf{h}^{(2)}_i+(1-\mathbf{g}_i)\odot\mathbf{s}_i .
\end{equation}
The head predicts the \emph{change} over persistence (A9), scaled by $\sigma_i$, the training-set
standard deviation of one-hour changes at station $i$, for each step $h=1..6$:
\begin{equation}
\hat y_{i,0.5}(\tau+h\Delta)=y_i(\tau^{\ast}_i)+\sigma_i\,\mathbf{w}_{h}^{\top}\mathbf{f}_i,\qquad
\hat y_{i,0.5\pm0.4}=\hat y_{i,0.5}\pm\sigma_i\,\mathrm{softplus}\big(\mathbf{w}^{\pm\top}_{h}\mathbf{f}_i\big),
\end{equation}
where $y_i(\tau^{\ast}_i)$ is the last valid level. The softplus offsets prevent quantile crossing.
The published forecast is $h=6$.

\textbf{(4) Training objective.}
\begin{equation}
\mathcal{L}=\frac{1}{|B|}\sum_{b\in B}\sum_{i\in\mathcal{S}}\sum_{h=1}^{6}\omega_{b}\,w_h\,M_{b,i,h}
\sum_{q\in\{0.1,0.5,0.9\}}\rho_q\!\left(\frac{y_i(\tau_b+h\Delta)-\hat y_{i,q}(\tau_b+h\Delta)}{\sigma_i}\right)
+\lambda\,\mathcal{R}_{\mathrm{mono}},
\end{equation}
with the pinball loss $\rho_q(r)=\max(qr,(q-1)r)$; $w_6=1$ and $w_{h<6}=0.2$ (auxiliary steps);
$M_{b,i,h}$ masking missing targets; $\omega_b=3$ for samples whose target exceeds the 95th
percentile level (floods matter most for preparedness) and 1 otherwise; and the optional monotonicity
penalty $\mathcal{R}_{\mathrm{mono}}=\mathbb{E}\big[\mathrm{ReLU}\big(-\partial\hat y_{i,0.5}/\partial Q^{\mathrm{sch}}\big)\big]$ (A9).
Training uses chronological splits (A1), AdamW, early stopping on the validation pinball loss, and the
missing/delay augmentation of Section~5.

\subsection*{How each driver reaches a forecast}
\begin{itemize}[leftmargin=*,itemsep=1pt]
\item \textbf{Upstream conditions.} S$k$ receives every upstream station through its own delay edge,
at lags matching the travel time. These edges carry most of the one-hour signal at S2--S4.
\item \textbf{Tributaries.} T1 is a node with lag windows into S2 (and onward to S3 and S4). The
unmonitored tributary is the latent node U. It is driven by rainfall, mostly R2 (A10), and is
implicitly calibrated by whatever part of S3's change S2 cannot explain, because this is the only
way the loss can use it.
\item \textbf{Rainfall.} Edges from R1 and R2 to their catchment nodes use lag windows covering
runoff times (0--6\,h). Rain therefore acts directly as local runoff and indirectly through the
upstream nodes and U.
\item \textbf{Reservoir releases.} Past releases act through the Res node. Announced releases act
through its future extension, and the lag window lets a release effective within the hour affect S1
before $\tau^{+}$ if its travel time is short enough (Section~5).
\end{itemize}

\subsection*{Physical assumptions versus learned mappings}
\textbf{Physical (fixed by design):} downstream edge directions (A7); lag windows from distances and
celerity bounds (A8); the residual-over-persistence form and the monotonicity prior (A9); the
existence of an unobserved inflow node between S2 and S3; and causality (no data after $\tau$, except
announced releases). \textbf{Learned:} travel time as a function of state (the lag weights $\alpha$), the
strength and attenuation of each influence ($W_V$, $\beta$, $\gamma$), the rainfall-runoff response,
the latent tributary state, the gating between own history and context, and the uncertainty
(quantiles). We do not hard-code a hydraulic model such as Muskingum or Saint-Venant, because its
parameters are not supplied; Section~6 describes how one can be added.

\subsection*{Justification of the main choices}
\begin{enumerate}[leftmargin=*,itemsep=1pt]
\item \textbf{Reach.} A forecaster cannot use information that is outside its reach, measured in hops
and in time (as the companion notebook shows). Transitive delay edges connect every upstream source
to every station within one layer, and the 12\,h window covers the longest admissible travel time.
A stack of single-hop layers would need at least four layers to bring reservoir information to S4.
\item \textbf{Bounded lag attention} is cheap, interpretable (the attention peak is an estimated
travel time, so the committee can inspect it), causal, and cannot use physically impossible lags.
Compared with full space-time attention it has few free parameters, which suits a few years of
data (A1).
\item \textbf{Separate encoders projected to a common width} handle the different sampling
(events versus grid) and feature sets of the sources (Figure~\ref{fig:arch}).
\item \textbf{Residual forecasting with gating} keeps persistence as a safe default. If upstream
inputs are missing, the gate falls back to the station's own history.
\item \textbf{Quantile outputs} give decision-makers an interval that can be read directly, for
example ``a 10\,\% chance that S3 exceeds the warning level within the hour''.
\end{enumerate}
